% Author's solution, supplied in this file on October 6, 2026.
% The author requested the geometric interpretation from section 5.1 in the body.
% Match the assignment's columns [a b]^T and [c d]^T, not the chapter's names.
\noindent\textbf{Solution.}
Let \((e_1,e_2)\) be the standard ordered basis of \(\R^2\). The column
vectors of the matrix are
\[
  v=ae_1+be_2,\qquad w=ce_1+de_2.
\]
Geometrically, \(v\) and \(w\) span the parallelogram
\[
  P(v,w)=\{sv+tw\mid 0\le s,t\le1\}.
\]
Let \(A\) be the oriented area function, normalized by \(A(e_1,e_2)=1\).
It is antisymmetric and bilinear. By uniqueness of the normalized
antisymmetric bilinear function, the determinant on columns satisfies
\[
  \det\begin{bmatrix}a&c\\b&d\end{bmatrix}=A(v,w).
\]
Thus the determinant is the signed area of this parallelogram, while its
ordinary area is \(|A(v,w)|\).

Bilinearity gives
\begin{align*}
  A(v,w)
    &=A(ae_1+be_2,ce_1+de_2)\\
    &=acA(e_1,e_1)+adA(e_1,e_2)\\
    &\quad+bcA(e_2,e_1)+bdA(e_2,e_2)\\
    &=adA(e_1,e_2)+bcA(e_2,e_1)\\
    &=(ad-bc)A(e_1,e_2)\\
    &=ad-bc.
\end{align*}
The terms with repeated basis vectors vanish because
\(A(e_1,e_1)=A(e_2,e_2)=0\). Antisymmetry gives
\(A(e_2,e_1)=-A(e_1,e_2)\), and the normalization gives the final equality.
Consequently, the determinant is \(ad-bc\). \qed
